% !TEX root = chapter-standalone.tex

\section{Canonical trace, duality}
\label{sec:canonical-trace-dual-objects}
\todotextjira{420}{\bernina: @JL: Seems like that here string diagrams would help.}

\

So far, a trace (or a feedback operator) on a symmetric strict monoidal category has been modeled as an operation of the following kind
\begin{equation*}
\generictracecat{X}{Y}{Z}:\HomSet{\stylecat{C}}{\styleobj{X}\mtimescat \styleobj{Z}}{\styleobj{Y}\mtimescat \styleobj{Z}}\to\HomSet{\stylecat{C}}{\styleobj{X}}{\styleobj{Y}}
\end{equation*}
defined for some, or all, morphisms of the above types. In terms of string diagrams, we have seen that we are able to draw this operation as a loop 
\begin{figure}[h!]
    \centering
    \begin{tikzpicture}[scale=0.7, every node/.style={transform shape},
        morbox/.style={draw=morphisms, fill=white, thick, minimum width=1.6cm, minimum height=2.2cm},
        wire/.style={objects, thick},
        wirelabel/.style={objects, font=\bfseries},
        tracebox/.style={draw=black, dashed, thick, fill=graphcolor, rounded corners=2pt},
    ]
        % Morphism box f
        \node[morbox] (f) at (2,0) {};
        \node[morphisms, font=\large\itshape] at (2,0) {$f$};

        % X wire: input from left
        \draw[wire] (-1.8, 0.55) -- (f.west |- 0,0.55);
        \node[wirelabel, above] at (-1.2, 0.55) {$\Obja$};

        % Y wire: output to right
        \draw[wire] (f.east |- 0,0.55) -- (5.8, 0.55);
        \node[wirelabel, above] at (5.2, 0.55) {$\Objb$};

        % Z wire: input side (from feedback loop up to f)
        \draw[wire] (0.6, -1.8) -- (0.6, -0.55) -- (f.west |- 0,-0.55);
        \node[wirelabel, left] at (0.6, -1.1) {$\Objc$};

        % Z wire: output side (from f down to feedback loop)
        \draw[wire] (f.east |- 0,-0.55) -- (3.4, -0.55) -- (3.4, -1.8);
        \node[wirelabel, right] at (3.4, -1.1) {$\Objc$};

        % Feedback loop: bottom horizontal segment
        \draw[wire] (0.6, -1.8) -- (3.4, -1.8);
    \end{tikzpicture}
    \caption{String diagram for the trace operator.}
\end{figure}


and incorporate it into the existing string diagram language with the help of the axioms in the definition of the trace. 


In the following we will get to know a special kind of symmetric strict monoidal category which is called compact closed, and in this case we will be able to add two new symbols to the string diagram language of symmetric strict monoidal categories. These new symbols represent special morphisms called `coevaluation', draw thus
\begin{figure}[h!]
    \centering
    \begin{tikzpicture}[scale=0.5, every node/.style={transform shape},
        wire/.style={objects, thick},
    ]
        % Coevaluation: C-shape opening to the right
        \draw[wire] (0, 1.2) -- (1.5, 1.2);
        \draw[wire] (0, 1.2) -- (0, -1.2);
        \draw[wire] (0, -1.2) -- (1.5, -1.2);
    \end{tikzpicture}
    \caption{Coevaluation.}
    \label{fig:coev-string}
\end{figure}

and `evaluation', drawn thus:

\begin{figure}[h!]
    \centering
    \begin{tikzpicture}[scale=0.5, every node/.style={transform shape},
        wire/.style={objects, thick},
    ]
        % Evaluation: reversed C-shape, opening to the left
        \draw[wire] (-1.5, 1.2) -- (0, 1.2);
        \draw[wire] (0, 1.2) -- (0, -1.2);
        \draw[wire] (-1.5, -1.2) -- (0, -1.2);
    \end{tikzpicture}
    \caption{Evaluation.}
    \label{fig:ev-string}
\end{figure}



With these symbols added to our string diagram language we will be able to `build' a trace in a canonical way from smaller building blocks, using only series and parallel compositions:

\begin{figure*}[h!]
    \centering
    \begin{tikzpicture}[
        morbox/.style={draw=morphisms, fill=white, thick, minimum width=1.2cm, minimum height=2.0cm},
        wire/.style={objects, thick},
        wirelabel/.style={font=\small, inner sep=2pt},
        boxlabel/.style={inner sep=1pt},
        shadebox/.style={fill=graphcolor, draw=graphcolorbord, thin, dashed, rounded corners=1pt},
    ]
        % Rows (y): 1.0 = top wire, -0.3 = middle wire, -1.45 = bottom wire;
        % 1.38 = labels of the top boxes, -1.88 = labels of the bottom boxes.

        % --- Left-hand side: the trace diagram ---
        \begin{scope}
            \draw[shadebox, rounded corners=2pt] (0.55,-1.6) rectangle (3.65,1.75);
            \node[morbox] (fL) at (2.1,0.35) {$\mora$};
            % X and Y wires
            \draw[wire] (-0.2,1.0) -- (fL.west |- 0,1.0);
            \draw[wire] (fL.east |- 0,1.0) -- (4.4,1.0);
            \node[wirelabel, above] at (0.15,1.0) {$\Obja$};
            \node[wirelabel, above] at (4.05,1.0) {$\Objb$};
            % Z feedback loop
            \draw[wire] (fL.east |- 0,-0.3) -- (3.15,-0.3) -- (3.15,-1.25) -- (1.05,-1.25) -- (1.05,-0.3) -- (fL.west |- 0,-0.3);
            \node[wirelabel, left] at (1.05,-0.8) {$\Objc$};
            \node[wirelabel, right] at (3.15,-0.8) {$\Objc$};
        \end{scope}

        % Equals sign
        \node[font=\Large] at (4.9,-0.1) {$=$};

        % --- Right-hand side: the canonical construction ---
        % Four slices, composed in series from left to right; in each slice
        % the shaded boxes are composed in parallel from top to bottom.
        \begin{scope}[xshift=5.4cm]
            % Shaded boxes (drawn first, behind everything)
            % Slice 1: id_X and coev_Z
            \draw[shadebox] (0.7,0.5) rectangle (2.2,1.75);
            \draw[shadebox] (0.7,-2.25) rectangle (2.2,0.2);
            % Slice 2: f and id on the dual of Z
            \draw[shadebox] (2.9,-0.9) rectangle (4.9,1.75);
            \draw[shadebox] (2.9,-2.25) rectangle (4.9,-1.1);
            % Slices 3 and 4, top: id_Y
            \draw[shadebox] (5.6,0.5) rectangle (9.4,1.75);
            % Slice 3, bottom: braiding
            \draw[shadebox] (5.6,-2.25) rectangle (7.2,0.2);
            % Slice 4, bottom: ev_Z
            \draw[shadebox] (7.9,-2.25) rectangle (9.4,0.2);

            % Morphism box f
            \node[morbox] (fR) at (3.9,0.35) {$\mora$};

            % Top wire: X into f, Y out of f
            \draw[wire] (0,1.0) -- (fR.west |- 0,1.0);
            \draw[wire] (fR.east |- 0,1.0) -- (10.1,1.0);

            % Coevaluation: C-shape; its upper end goes into f,
            % its lower end passes below f to the braiding
            \draw[wire] (fR.west |- 0,-0.3) -- (1.2,-0.3) -- (1.2,-1.45) -- (5.9,-1.45);

            % Z out of f to the braiding
            \draw[wire] (fR.east |- 0,-0.3) -- (5.9,-0.3);

            % Braiding: straight X-crossing, no gap (symmetric case)
            \draw[wire] (5.9,-0.3) -- (6.9,-1.45);
            \draw[wire] (5.9,-1.45) -- (6.9,-0.3);

            % Evaluation: reversed C-shape
            \draw[wire] (6.9,-0.3) -- (8.9,-0.3) -- (8.9,-1.45) -- (6.9,-1.45);

            % Labels of the boxes
            \node[boxlabel] at (1.45,1.38) {$\catidat{\Obja}$};
            \node[boxlabel] at (7.5,1.38) {$\catidat{\Objb}$};
            \node[boxlabel] at (1.45,-1.88) {$\coev_{\Objc}$};
            \node[boxlabel] at (3.9,-1.88) {$\catidat{\dualof{\Objc}}$};
            \node[boxlabel] at (6.4,-1.88) {$\braiding_{\Objc,\dualof{\Objc}}$};
            \node[boxlabel] at (8.65,-1.88) {$\ev_{\Objc}$};

            % Labels of the wires
            \node[wirelabel, above] at (0.35,1.0) {$\Obja$};
            \node[wirelabel, above] at (2.55,1.0) {$\Obja$};
            \node[wirelabel, above] at (5.25,1.0) {$\Objb$};
            \node[wirelabel, above] at (9.75,1.0) {$\Objb$};
            \node[wirelabel, above] at (2.55,-0.3) {$\Objc$};
            \node[wirelabel, above] at (5.25,-0.3) {$\Objc$};
            \node[wirelabel, above] at (2.55,-1.45) {$\dualof{\Objc}$};
            \node[wirelabel, above] at (5.25,-1.45) {$\dualof{\Objc}$};
            \node[wirelabel, above] at (7.55,-0.3) {$\dualof{\Objc}$};
            \node[wirelabel, above] at (7.55,-1.45) {$\Objc$};
        \end{scope}
    \end{tikzpicture}
    \caption{The canonical trace, built from coevaluation, the morphism~$f$, braiding, and evaluation.}
    \label{fig:canonical-trace-construction}
\end{figure*}


In particular, when given a compact closed category, there is a theorem that guarantees for us that this construction gives a trace, and we then don't need to go and check the many axioms in the definition of trace. 

So what is a compact closed category? As a first slogan, it is a symmetric strict monoidal category where each object has a `dual object'. 

Next you may wonder what `dual objects' are. Well, glad you asked! Let's have a look. 


\



\

\subsection{Definitions}

\begin{ctdefinition}[Dual object]
    \label{def:dualizable-object}
    Let~$\tup{\CatC,\mtimesC,\idmoncat_{\CatC}}$ be a \SY{strict monoidal category}, and let~$\Obja \setin \ObC$.
    A \maindef{right dual object} of~$\Obja$ is specified by:

    \constit
    \begin{enumerate}
        \item an object $\Objadual \setin \ObC$;
        \item a morphism $\ev_\Obja\colon \Objadual \mtimescatob \Obja \mto \idmoncat$, called \emph{evaluation}; \label{def:ev1}
        \item a morphism $\coev_\Obja\colon \idmoncat \mto \Obja \mtimescatob \Objadual$, called \emph{coevaluation};\label{def:coev1}
    \end{enumerate}

    \condit
    \begin{enumerate}
        \item
              \begin{equation}\label{eq:snake-eq-1}
                  (\coev_\Obja \mtimescatmor \catidat\Obja) \mthen (\catidat\Obja \mtimescatmor \ev_\Obja) = \catidat\Obja;
              \end{equation}
        \item
              \begin{equation}\label{eq:snack-eq-2}
                  (\catidat{\Objadual} \mtimescatmor \coev_\Obja) \mthen (\ev_\Obja \mtimescatmor \catidat{\Objadual}) = \catidat{\Objadual}.
              \end{equation}
    \end{enumerate}
\end{ctdefinition}

These compositions make sense because the monoidal category is strict.
For example, in \cref{eq:snake-eq-1}, the source of $\coev_\Obja \mtimescatmor \catidat\Obja$ is $\idmoncat \mtimescatob \Obja = \Obja$, its target is $\Obja \mtimescatob \Objadual \mtimescatob \Obja$, which is also the source of $\catidat\Obja \mtimescatmor \ev_\Obja$, and the target of the latter is $\Obja \mtimescatob \idmoncat = \Obja$.

The two conditions above are called the \emph{snake equations}.
In terms of string diagrams, they say that the following two compositions each equal the identity:

\begin{figure}[h!]
    \centering
    \begin{tikzpicture}[scale=0.6, every node/.style={transform shape},
        wire/.style={objects, thick},
        wirelabel/.style={objects, font=\small},
    ]
        % --- Snake equation 1: S-shape ---
        % Three levels: top (1.5), middle (0), bottom (-1.5)
        % Coev on left connects top and middle; ev on right connects middle and bottom.
        % Input X at bottom-left, output X at top-right.

        % Output X exits at top-right
        \draw[wire] (0, 1.5) -- (5.0, 1.5);
        \node[wirelabel, above] at (4.7, 1.5) {$\Obja$};

        % Coev: vertical on left from top to middle
        \draw[wire] (0, 1.5) -- (0, 0);
        \node[wirelabel, left] at (-0.1, 0.75) {\scriptsize coev};

        % Horizontal middle wire from coev to ev
        \draw[wire] (0, 0) -- (3.5, 0);

        % Ev: vertical on right from middle to bottom
        \draw[wire] (3.5, 0) -- (3.5, -1.5);
        \node[wirelabel, right] at (3.6, -0.75) {\scriptsize ev};

        % Input X enters at bottom-left
        \draw[wire] (-1.5, -1.5) -- (3.5, -1.5);
        \node[wirelabel, above] at (-1.2, -1.5) {$\Obja$};

        % Equals sign
        \node[font=\Large] at (6.5, 0) {$=$};

        % Straight X wire (identity)
        \draw[wire] (8.0, 0) -- (10.5, 0);
        \node[wirelabel, above] at (8.3, 0) {$\Obja$};
        \node[wirelabel, above] at (10.2, 0) {$\Obja$};
    \end{tikzpicture}
    \caption{Snake equation 1.}
    \label{fig:snake-eq-1-string}
\end{figure}

\begin{figure}[h!]
    \centering
    \begin{tikzpicture}[scale=0.6, every node/.style={transform shape},
        wire/.style={objects, thick},
        wirelabel/.style={objects, font=\small},
    ]
        % --- Snake equation 2: mirrored S-shape ---
        % Three levels: top (1.5), middle (0), bottom (-1.5)
        % Ev on RIGHT connects top and middle; coev on LEFT connects middle and bottom.
        % Input X-dual at top-left, output X-dual at bottom-right.

        % Input X-dual enters at top-left, goes right to ev
        \draw[wire] (-1.5, 1.5) -- (3.5, 1.5);
        \node[wirelabel, above] at (-1.2, 1.5) {$\Objadual$};

        % Ev: vertical on RIGHT from top to middle
        \draw[wire] (3.5, 1.5) -- (3.5, 0);
        \node[wirelabel, right] at (3.6, 0.75) {\scriptsize ev};

        % Horizontal middle wire from ev back left to coev
        \draw[wire] (0, 0) -- (3.5, 0);

        % Coev: vertical on LEFT from middle to bottom
        \draw[wire] (0, 0) -- (0, -1.5);
        \node[wirelabel, left] at (-0.1, -0.75) {\scriptsize coev};

        % Output X-dual exits at bottom-right
        \draw[wire] (0, -1.5) -- (5.0, -1.5);
        \node[wirelabel, above] at (4.7, -1.5) {$\Objadual$};

        % Equals sign
        \node[font=\Large] at (6.5, 0) {$=$};

        % Straight X-dual wire (identity)
        \draw[wire] (8.0, 0) -- (10.5, 0);
        \node[wirelabel, above] at (8.3, 0) {$\Objadual$};
        \node[wirelabel, above] at (10.2, 0) {$\Objadual$};
    \end{tikzpicture}
    \caption{Snake equation 2.}
    \label{fig:snake-eq-2-string}
\end{figure}

\begin{ctdefinition}\label{def:right-dualizable}
    An object $\Obja$ in a \SY{strict monoidal category} is called \maindef{right dualizable} if there exists, in the category, a right dual object of $\Obja$.
\end{ctdefinition}

\begin{remark}\label{rem:left-dualizability}
    There is an analogous definition of \maindef{left dual object} and \maindef{left dualizability}.
    One can show that when the \SY{strict monoidal category} in question is symmetric, then left dual objects can be seen as right dual objects, and vice versa.
    In this case, we then speak simply of \maindef{dualizability}.
\end{remark}

\begin{ctdefinition}[Compact closed category]
    \label{def:compact-closed-category}
    A symmetric \SY{strict monoidal category} (\cref{def:sym-stacking-cat}) is called \emph{compact closed} if every object is dualizable.
\end{ctdefinition}


\begin{example}
    The symmetric strict monoidal category $\RelL$ of relations between sets of tuples (\cref{exa:RelL-symmetric}) is compact closed.
    Every object is its own dual; the evaluation and coevaluation morphisms are described in \cref{ex:RelDualsTrace}.
\end{example}

\begin{example}
    \label{exa:mRelL-compact-closed}
    The symmetric strict monoidal category $\mRelL$ of monotone relations between preorders of tuples (\cref{exa:RelL-symmetric}) is compact closed.
    The dual of an object $\posA$ is its opposite preorder $\posA\op$, in which the order is reversed.
    Since tuples are compared entry by entry, this is again a preorder of tuples.
    The evaluation and the coevaluation are the monotone relations
    \begin{equation}
        \label{eq:mRelL-ev}
        \ev_{\posA} \colon \posA\op \mtimescatob \posA \profto \emptycobj,
        \qquad
        \ev_{\posA} = \makeset{ \tup{\maketupconcat{\posela, \posela'}, \emptytuple} \mid \posela' \posAleq \posela }
    \end{equation}
    and
    \begin{equation}
        \label{eq:mRelL-coev}
        \coev_{\posA} \colon \emptycobj \profto \posA \mtimescatob \posA\op,
        \qquad
        \coev_{\posA} = \makeset{ \tup{\emptytuple, \maketupconcat{\posela, \posela'}} \mid \posela' \posAleq \posela },
    \end{equation}
    where $\posela$ and $\posela'$ are elements of the underlying set of $\posA$, which is also the underlying set of $\posA\op$.
    Compared to $\RelL$, the equality of the two elements is replaced by an inequality.
    Written as Boolean profunctors, the same data show that $\ProfBoolL$ is compact closed, and they restrict to the full subcategories $\feasL$ and $\DPL$ of \cref{exa:codesign-strict-monoidal}; see \cref{ex:DPSnakeTrace}.
\end{example}

\begin{example}[Categories which are not compact closed]
    \label{exa:not-compact-closed}
    Not every symmetric strict monoidal category is compact closed.
    Here are three examples in which very few objects are dualizable.
    In each of them, we look at the left-hand side of the first snake equation \cref{eq:snake-eq-1}.
    \begin{itemize}
        \item In the category $\SetL$ of functions between sets of tuples (\cref{exa:SetL-symmetric}), the monoidal unit $\emptycobj$ has a single element.
              So for an object $\setA$ with a right dual $\setA\dual$, the coevaluation $\emptycobj \mto \setA \cprod \setA\dual$ picks out one element $\maketupconcat{\ela_0, \ela_0'}$, and the evaluation is the only function $\setA\dual \cprod \setA \mto \emptycobj$.
              The left-hand side of \cref{eq:snake-eq-1} therefore maps an element $\ela$ of $\setA$ first to $\maketupconcat{\ela_0, \ela_0', \ela}$, and then to $\ela_0$.
              This is a constant function, and it is the identity function only if $\ela_0$ is the only element of $\setA$.
              So in $\SetL$, only the sets with exactly one element are dualizable.
        \item In the category $\MatR$ of real matrices with the monoidal product which adds dimensions (\cref{exa:MatR-strict-monoidal}), the monoidal unit is $\reals^0$.
              So for an object $\reals^n$ with a right dual $\reals^m$, the coevaluation $\reals^0 \mto \reals^{n+m}$ is the matrix without columns, and the evaluation $\reals^{m+n} \mto \reals^0$ is the matrix without rows.
              The left-hand side of \cref{eq:snake-eq-1} therefore maps a vector $\vectob{v} \setin \reals^n$ first to the vector in $\reals^{n+m+n}$ whose last $n$ entries are those of $\vectob{v}$, and whose other entries are zero, and then to the first $n$ entries of that vector, which are zero.
              This is the zero matrix, and it is the identity matrix only if $n = 0$.
              So here only the monoidal unit is dualizable.
              This fits with what we saw in \cref{sec:feedback-MatR}: the formula for closing loops is not defined for all matrices.
        \item In the category $\RelR$ of relations between sets of indexed elements (\cref{exa:SetR-symmetric}), the monoidal unit is the empty set.
              So for an object $\setA$ with a right dual $\setA\dual$, both the coevaluation $\Emptyset \mto \setA \sprod \setA\dual$ and the evaluation $\setA\dual \sprod \setA \mto \Emptyset$ are empty relations, and the left-hand side of \cref{eq:snake-eq-1} is the empty relation as well.
              It is the identity relation only if $\setA$ is empty.
              So here, too, only the monoidal unit is dualizable.
              Nevertheless, $\RelR$ has a trace (\cref{def:RelR-feedback}).
              This shows that a trace does not have to come from dual objects in the way which we describe in this section.
    \end{itemize}
\end{example}


\begin{ctdefinition}[Dual morphism]\label{def:dual-morphism}
    Let $\tup{\CatC,\mtimesC,\idmoncat_{\CatC}}$ be a \SY{strict monoidal category}, let $\mora \colon \Obja \mto \Objb$ be a morphism in it, and suppose that $\Obja$ and $\Objb$ have right duals $\Obja\dual$ and $\Objb\dual$, respectively.
    The \emph{right dual} of $\mora$ is the morphism $\mora\dual \colon \Objb\dual \mto \Obja\dual$ given by the composition
    \begin{equation}
        \label{eq:dual_morphism}
        \mora\dual \definedas
        (\catidat{\Objb\dual} \mtimescatmor \coev_\Obja)
        \mthen (\catidat{\Objb\dual} \mtimescatmor \mora \mtimescatmor \catidat{\Obja\dual})
        \mthen (\ev_\Objb \mtimescatmor \catidat{\Obja\dual}).
    \end{equation}
\end{ctdefinition}

In terms of string diagrams, the dual morphism is illustrated as follows:

\begin{figure}[h!]
    \centering
    \begin{tikzpicture}[scale=0.6, every node/.style={transform shape},
        morbox/.style={draw=morphisms, fill=white, thick, minimum width=1.2cm, minimum height=1.2cm},
        wire/.style={objects, thick},
        wirelabel/.style={objects, font=\small},
    ]
        % Backward S-shape: Y-dual enters top-left, X-dual exits bottom-right.
        % Three levels: top (1.5), middle (0), bottom (-1.5)
        % Ev on right connects top to middle; coev on left connects middle to bottom.
        % f centered at middle level between coev and ev.

        % Y-dual enters at top-left, goes right to ev
        \draw[wire] (-1.5, 1.5) -- (5.5, 1.5);
        \node[wirelabel, above] at (-1.2, 1.5) {$\Objb\dual$};

        % Ev: vertical on right from top to middle
        \draw[wire] (5.5, 1.5) -- (5.5, 0);
        \node[wirelabel, right, font=\scriptsize] at (5.6, 0.75) {ev};

        % f box at middle level, centered between coev and ev
        \node[morbox] (f) at (3.0, 0) {};
        \node[morphisms, font=\itshape] at (3.0, 0) {$\mora$};

        % Y wire from ev goes left into f
        \draw[wire] (5.5, 0) -- (f.east |- 0,0);
        \node[wirelabel, above] at (4.5, 0) {$\Objb$};

        % X wire from f goes left to coev
        \draw[wire] (f.west |- 0,0) -- (0.5, 0);
        \node[wirelabel, above] at (1.5, 0) {$\Obja$};

        % Coev: vertical on left from middle to bottom
        \draw[wire] (0.5, 0) -- (0.5, -1.5);
        \node[wirelabel, left, font=\scriptsize] at (0.4, -0.75) {coev};

        % X-dual exits at bottom-right
        \draw[wire] (0.5, -1.5) -- (7.0, -1.5);
        \node[wirelabel, above] at (6.7, -1.5) {$\Obja\dual$};
    \end{tikzpicture}
    \caption{String diagram for the dual morphism.}
    \label{fig:dual-morphism-string}
\end{figure}

\begin{gradedexercise}[\exname{RelDualsTrace}]
    \label{ex:RelDualsTrace}

    In this exercise we work with the symmetric strict monoidal category $\RelL$ of relations between sets of tuples (\cref{def:RelL,exa:RelL-symmetric}).
    Recall that its monoidal unit is the set $\emptycobj$, whose only element is the empty tuple $\emptytuple$, and that its braiding is
    \begin{equation}
        \braiding_{\setA, \setB} = \makeset{ \tup{\maketupconcat{\ela, \elb}, \maketupconcat{\elb, \ela}} \mid \ela \setin \setA \text{ and } \elb \setin \setB }.
    \end{equation}

    This symmetric strict monoidal category is compact closed when we let the dual $\setA \dual$ of any object $\setA$ be the object itself, $\setA \dual = \setA$, and we define evaluation and coevaluation by
    \begin{equation}
        \ev_\setA \colon \setA \cprod \setA \mto \emptycobj, \quad \ev_\setA = \makeset{\tup{\maketupconcat{\ela, \elb}, \emptytuple} \mid \ela, \elb \setin \setA \text{ and } \ela = \elb}
    \end{equation}
    and
    \begin{equation}
        \coev_\setA \colon \emptycobj \mto \setA \cprod \setA, \quad \coev_\setA = \makeset{\tup{\emptytuple, \maketupconcat{\ela, \elb}} \mid \ela, \elb \setin \setA \text{ and } \ela = \elb}
    \end{equation}
    respectively.

    Your tasks:
    \begin{enumerate}
        \item
              Let $\relB \colon \setA \mto \setB$ be a (generic) morphism in $\RelL$.
              Compute the dual morphism $\relB\dual \colon \setB\dual \mto \setA\dual$ (\cref{def:dual-morphism}).

        \item
              Let $\relA \colon \setA \cprod \setC \mto \setB \cprod \setC$ be a morphism in $\RelL$.
              Show that the composition
              \begin{equation}
                  (\catid_\setA \mtimescatmor \coev_\setC) \mthen (\relA \mtimescatmor \catid_\setC) \mthen (\catid_\setB \mtimescatmor \braiding_{\setC, \setC}) \mthen (\catid_\setB \mtimescatmor \ev_\setC)
              \end{equation}
              is equal to the relation
              \begin{equation}
                  \makeset{\tup{\ela, \elb} \setin\setA \cartprod \setB \mid \exists \ \elc \setin \setC \colon \tup{\maketupconcat{\ela, \elc},\maketupconcat{\elb,\elc}} \setin \relA},
              \end{equation}
              that is, to the trace of $\relA$ as defined in \cref{def:RelL-feedback}.
    \end{enumerate}
\end{gradedexercise}
\solutionof{RelDualsTrace}



\subsection{Why the names ``dual" and ``(co)evaluation''?}

The names come from linear algebra, and more precisely from the tensor product of vector spaces.
Recall from \cref{exa:MatR-Kronecker-strict-monoidal} that the Kronecker product of matrices is the tensor product, written in coordinates.
We first describe the situation for vector spaces in general, without coordinates, and then translate it into our strict monoidal category of matrices.

Given a \SY{vector space} $\styleobj{V} $, its \emph{linear dual} is the real vector space
\begin{equation}
    \label{eq:lin-dual-vecsp}
    \styleobj{V} ^{*} := \makeset{ \text{linear maps } \styleobj{V} \mto \reals }.
\end{equation}
%
Recall from linear algebra the following fact about any \SY{vector space} $\styleobj{V} $:
%
\begin{equation}
    \label{eq:double-dual-fin-dim}
    \styleobj{V} \simeq (\styleobj{V} ^{*})^* \text{ if and only if } \dim \styleobj{V} < \infty.
\end{equation}
%
Thus we may say that the finite-dimensional real \SY{vector spaces} are characterizable based on their behavior in this way with respect to the operation of taking the linear dual. 


In the following, $\styleobj{V} $ denotes a \emph{finite-dimensional} real \SY{vector space}.
The evaluation map $\ev_\styleobj{V} $ associated to $\styleobj{V} $ is
\begin{equation}
    \label{eq:evaluation-vecsp}
    \defmapperiod{
        \ev_\styleobj{V}
    }{
        \styleobj{V}^* \otimes \styleobj{V}
    }{
        \mto
    }{
        \reals
    }{
        \tup{l, v}
    }{
        l(v)
    }
\end{equation}
In other words, given $\tup{l, v}$, the map $\ev_\styleobj{V} $ evaluates $l$ at $v$.

The coevaluation map $\coev_\styleobj{V} $ associated to $\styleobj{V} $ is slightly trickier to describe.
Let $\makeset{e_1, \dots, e_n }$ be a basis of $\styleobj{V}$, and let $\makeset{e_1^*, \dots, e_n^* }$ be the corresponding dual basis of~$\styleobj{V}^*$.
Then
\begin{equation}
    \label{eq:coevaluation-vecsp}
    \defmapperiod{
        \coev_\styleobj{V}
    }{
        \reals
    }{
        \mto
    }{
        \styleobj{V} \otimes \styleobj{V}^*
    }{
        \lambda
    }{
        \lambda \sum_{i=1}^n e_i \otimes e_i^*
    }
\end{equation}
It turns out that this map is independent of the choice of basis.
One way to think of this coevaluation map is to recall that $\styleobj{V} \otimes \styleobj{V}^* \simeq \Hom(\styleobj{V}, \styleobj{V})$.
Under this identification, $\coev_\styleobj{V}$ maps any scalar $\lambda$ to the linear \SY{endomorphism} of $\styleobj{V}$ which is ``multiplication by $\lambda$''.
(In terms of matrices, this is a diagonal matrix, with $\lambda$ at every entry of the diagonal.)

Let us now translate this into our strict monoidal category $\MatR$ of real matrices with the Kronecker product (\cref{exa:MatR-Kronecker-strict-monoidal,exa:MatR-Kronecker-symmetric}), in which the vector space $\styleobj{V}$ becomes a coordinate space $\reals^n$.
A linear map $\reals^n \mto \reals$ is a matrix with one row and $n$ columns, and by transposing it we obtain a matrix with $n$ rows and one column, that is, an element of $\reals^n$.
In this sense the linear dual of $\reals^n$ is $\reals^n$ again, and so we let every object be its own dual:
\begin{equation}
    \dualof{(\reals^n)} \definedas \reals^n.
\end{equation}
The evaluation map \cref{eq:evaluation-vecsp} becomes the matrix
\begin{equation}
    \label{eq:evaluation-MatR}
    \ev_{\reals^n} \colon \reals^n \mtimescat \reals^n \mto \reals^1
    \qquad \text{with} \qquad
    \ev_{\reals^n}(\vectob{l} \mtimescat \vectob{v}) = \sum_{i=1}^n l_i v_i
\end{equation}
for all $\vectob{l}, \vectob{v} \setin \reals^n$; it evaluates the linear map with the coefficients $l_1, \ldots, l_n$ at the vector $\vectob{v}$.
The coevaluation map \cref{eq:coevaluation-vecsp} becomes the matrix
\begin{equation}
    \label{eq:coevaluation-MatR}
    \coev_{\reals^n} \colon \reals^1 \mto \reals^n \mtimescat \reals^n
    \qquad \text{with} \qquad
    \coev_{\reals^n}(\lambda) = \lambda \sum_{i=1}^n \vectob{e}_i \mtimescat \vectob{e}_i,
\end{equation}
where $\vectob{e}_1, \ldots, \vectob{e}_n$ is the standard basis of $\reals^n$.
Concretely, $\ev_{\reals^n}$ is the matrix with one row and $n^2$ columns, and $\coev_{\reals^n}$ is the matrix with $n^2$ rows and one column, whose entries are the entries of the identity matrix $\idmat_n$, listed row by row.

\begin{proposition}
    \label{prop:MatR-Kronecker-compact-closed}
    Consider the symmetric strict monoidal category $\MatR$ of real matrices with the Kronecker product as monoidal product (\cref{exa:MatR-Kronecker-symmetric}).
    With the evaluation and coevaluation morphisms \cref{eq:evaluation-MatR} and \cref{eq:coevaluation-MatR}, every object $\reals^n$ is its own dual, and so this symmetric strict monoidal category is compact closed.
\end{proposition}


\begin{gradedexercise}[\exname{VectSnakeEquations}]
    \label{ex:VectSnakeEquations}
    Check the snake equations (\cref{eq:snake-eq-1,eq:snack-eq-2}) for the strict monoidal category $\MatR$ of real matrices with the Kronecker product, where the evaluation and coevaluation morphisms are the matrices \cref{eq:evaluation-MatR} and \cref{eq:coevaluation-MatR}.
\end{gradedexercise}

\solutionof{VectSnakeEquations}

\begin{example}[The dual of a matrix is its transpose]
    \label{exa:MatR-dual-morphism}
    Let us compute dual morphisms (\cref{def:dual-morphism}) in the category $\MatR$ of real matrices with the Kronecker product, with the duality data of \cref{prop:MatR-Kronecker-compact-closed}.
    Let $\matmor{A} \colon \reals^k \mto \reals^m$ be a matrix with entries $a_{ij}$.
    Its dual is a matrix $\matmor{A}\dual \colon \reals^m \mto \reals^k$, and by \cref{eq:dual_morphism} it maps a vector $\vectob{w} \setin \reals^m$ as follows, where $\vectob{e}_1, \ldots, \vectob{e}_k$ is the standard basis of $\reals^k$:
    \begin{equation}
        \vectob{w}
        \ \longmapsto \
        \sum_{j=1}^k \vectob{w} \mtimescat \vectob{e}_j \mtimescat \vectob{e}_j
        \ \longmapsto \
        \sum_{j=1}^k \vectob{w} \mtimescat \matmor{A}\vectob{e}_j \mtimescat \vectob{e}_j
        \ \longmapsto \
        \sum_{j=1}^k \Bigl( \sum_{i=1}^m w_i a_{ij} \Bigr) \vectob{e}_j.
    \end{equation}
    In the last step we used that $\matmor{A}\vectob{e}_j$ is the column $j$ of $\matmor{A}$, and that the evaluation \cref{eq:evaluation-MatR} multiplies the entries of $\vectob{w}$ with the entries of this column and adds up the results.
    So the entry $j$ of $\matmor{A}\dual\vectob{w}$ is $\sum_{i=1}^m a_{ij} w_i$, which means that $\matmor{A}\dual$ is the \emph{transpose} of the matrix $\matmor{A}$.
    In \cref{ex:RelDualsTrace} you are asked to find out what the dual of a relation is.
\end{example}

%The equations \cref{eq:dualizability-cond-1} and \cref{eq:dualizability-cond-2} form the basis for the general notion of dualizability in a \SY{monoidal category}.


\todotext{J: Add remark explaining that Vect has duality data only for finite-dimensional vector spaces, even though definition of 'linear dual' works also for finite dimenions. Also we can maybe then make a general remark somewhere that in any CCC the dual of an object X has X itself as its dual and X and its double dual are always isomorphic. Also maybe clarify somewhere the question of existence and uniqueness with respect to definition of CCC.... we just ask for dualizability but not duality data? And how unique is the latter? I think this is currently actually clearly treated in the text, but maybe warrants a remark and some additional further material..}


\subsection{Canonical trace}
\linkvideo{spring2021-par-feedback:feedback} % Feedback
\todotextjira{199}{\alphubel: @JL: Write something about how we already saw string diagrams, and that (co)evaluation maps allow for bending the wires around to change the their direction, and that we can encode dual objects by adding a ``flow'' direction to the wires.
    Then motivate the notion of trace in a \SY{monoidal category} by asking what we get when we bend a wire around and let the outout of a morphism flow back and be it's input -- ``closing the loop''}

%\linkvideo{spring2021-par-feedback:mon-cat:string-diag} % String diagrams

\begin{ctdefinition}[Trace of an endomorphism]
    \label{def:trace_endo}
    Let~$\tup{\CatC, \mtimescat, \idmoncat, \braiding}$ be a symmetric \SY{strict monoidal category}.
    Let~$\Obja \setin \ObC$ be dualizable and let
    \begin{equation}
        \mora \colon \Obja \mto \Obja.
    \end{equation}
    The \emph{trace} of~$\mora$ is the morphism
    \begin{equation}
        \Tr(\mora) \colon \idmoncat \mto \idmoncat
    \end{equation}
    defined by
    \begin{equation}
        \label{eq:endo_trace}
        \Tr(\mora) \definedas \coev_\Obja \mthen (\mora \mtimescatmor \catidat{\Objadual}) \mthen \braiding_{\Obja,\Objadual} \mthen \ev_\Obja.
    \end{equation}
\end{ctdefinition}

\begin{gradedexercise}[\exname{LinearAlgebraTrace}]
    \label{ex:LinearAlgebraTrace}
    Let~\CatC be the symmetric strict monoidal category $\MatR$ of real matrices with the Kronecker product as monoidal product (\cref{exa:MatR-Kronecker-symmetric}).
    We saw in \cref{prop:MatR-Kronecker-compact-closed} that every object $\reals^n$ of this category is dualizable: it is its own dual, with the evaluation and coevaluation matrices
    \begin{equation}
        \ev_{\reals^n} \colon \reals^n \mtimescat \reals^n \mto \reals^1,
        \qquad
        \ev_{\reals^n}(\vectob{l} \mtimescat \vectob{v}) = \sum_{i=1}^n l_i v_i,
    \end{equation}
    and
    \begin{equation}
        \coev_{\reals^n} \colon \reals^1 \mto \reals^n \mtimescat \reals^n,
        \qquad
        \coev_{\reals^n}(\lambda) = \lambda \sum_{i=1}^n \vectob{e}_i \mtimescat \vectob{e}_i,
    \end{equation}
    where $\vectob{e}_1, \ldots, \vectob{e}_n$ is the standard basis of $\reals^n$.

    Let~$\mat{M}$ be a square matrix with $n$ rows and $n$ columns, that is, a morphism $\reals^n \mto \reals^n$.
    Compute its trace~$\Tr(\mat{M}) \colon \reals^1 \mto \reals^1$ according to \cref{def:trace_endo}, and explain why it is the $1 \times 1$ matrix whose entry is the trace of~$\mat{M}$, where ``trace'' in this latter phrase is the usual notion that we know from linear algebra.
\end{gradedexercise}
\solutionof{LinearAlgebraTrace}

\begin{gradedexercise}[\exname{DPSnakeTrace}]
    \label{ex:DPSnakeTrace}

    In this exercise we work with the symmetric strict monoidal category $\DPL$ of design problems between posets of tuples (\cref{exa:codesign-strict-monoidal,exa:RelL-symmetric}).
    Recall that its monoidal unit is the poset $\emptycobj$, whose only element is the empty tuple $\emptytuple$, and that the braiding
    \begin{equation}
        \braiding_{\posgenA,\posgenB}\colon (\posgenA \mtimescatob \posgenB)\op \Ptimes (\posgenB \mtimescatob \posgenA) \mto \Bool
    \end{equation}
    is given by
    \begin{equation}
        \braiding_{\posgenA,\posgenB}(\maketupconcat{\ela_1, \elb_1}, \maketupconcat{\elb_2, \ela_2}) \definedas \pars{\ela_1\posleqof\posgenA \ela_2}\booland \pars{\elb_1\posleqof\posgenB \elb_2}.
    \end{equation}
    In the following you may use that
    \begin{equation}
        (\posgenA \mtimescatob \posgenB)\op = \posgenA\op \mtimescatob \posgenB\op
    \end{equation}
    for any objects $\posgenA$, $\posgenB$, because tuples are compared entry by entry.
    Also, recall that $(\posgenA\op)\op = \posgenA$.

    Let us define the following duality data:
    \begin{itemize}
        \item $\posgenA\dual \definedas \posgenA\op$
        \item $\ev_{\posgenA} \colon (\posgenA\op \mtimescatob \posgenA)\op \Ptimes \emptycobj \mto \Bool, \quad \tup{\maketupconcat{\ela, \elb}, \emptytuple} \mapsto \elb \posleqof\posgenA \ela$
        \item $\coev_{\posgenA} \colon \emptycobj\op \Ptimes (\posgenA \mtimescatob \posgenA\op) \mto \Bool, \quad \tup{\emptytuple, \maketupconcat{\ela, \elb}} \mapsto \elb \posleqof\posgenA \ela$
    \end{itemize}
    Here $\ela$ and $\elb$ are elements of the underlying set of $\posgenA$, which is also the underlying set of $\posgenA\op$.

    Your tasks:

    \begin{enumerate}
        \item
              Check that $\ev_{\posgenA}$ and $\coev_{\posgenA}$, as defined above, are morphisms in $\DPL$.
        \item
              For an arbitrary object $\posgenA$ and the duality data given above, prove that the snake equation
              \begin{equation}
                  (\coev_{\posgenA} \mtimescatmor \catidat \posgenA) \mthen (\catidat\posgenA \mtimescatmor \ev_{\posgenA}) = \catidat\posgenA
              \end{equation}
              holds.
    \end{enumerate}
\end{gradedexercise}
\solutionof{DPSnakeTrace}

\begin{ctdefinition}[Trace of a generalized endomorphism]
    \label{def:trace_gen_endo}
    Let~$\tup{\CatC, \mtimescat, \idmoncat, \braiding}$ be a symmetric \SY{strict monoidal category}.
    Let~$\Objc \setin \ObC$ be dualizable and let
    \begin{equation}
        \mora \colon \Obja \mtimescatob \Objc \mto \Objb \mtimescatob \Objc.
    \end{equation}
    The \emph{trace over}~$\Objc$ of~$\mora$ is the morphism
    \begin{equation}
        \Tr_{\Obja, \Objb}^\Objc (\mora) \colon \Obja \mto \Objb
    \end{equation}
    defined by
    \begin{equation}
        \label{eq:endo_trace_gen}
        \begin{aligned}
            \Tr_{\Obja, \Objb}^\Objc (\mora) \definedas {}
             & (\catidat\Obja \mtimescatmor \coev_\Objc)
            \mthen (\mora \mtimescatmor \catidat{\dualof{\Objc}}) \\
             & \mthen (\catidat\Objb \mtimescatmor \braiding_{\Objc,\dualof{\Objc}})
            \mthen (\catidat\Objb \mtimescatmor \ev_\Objc).
        \end{aligned}
    \end{equation}
\end{ctdefinition}

This is the construction which is shown in \cref{fig:canonical-trace-construction}; the trace of an endomorphism in \cref{def:trace_endo} is the special case in which $\Obja$ and $\Objb$ are the monoidal unit.
One can show that in a compact closed category, where every object is dualizable, these operations satisfy all the conditions of \cref{def:feedback-cat} and of \cref{def:trace-operator}.
The trace which is obtained in this way is called the \emph{canonical trace} of the compact closed category.
For $\RelL$, it is the trace of \cref{def:RelL-feedback} (\cref{ex:RelDualsTrace}), and for matrices with the Kronecker product, it is the trace of \cref{eq:MatR-Kronecker-trace}, which replaces each block of a matrix by its usual trace.

\begin{example}[Canonical trace for monotone relations]
    \label{exa:mRelL-canonical-trace}
    Let us compute the canonical trace of the compact closed category $\mRelL$, with the duality data of \cref{exa:mRelL-compact-closed}.
    Let $\monrelM \colon \posA \mtimescatob \posC \profto \posB \mtimescatob \posC$ be a monotone relation.
    We go through the four factors of \cref{eq:endo_trace_gen}.
    The coevaluation $\coev_{\posC}$ provides an element $\poselc$ of $\posC$ and an element $\poselc''$ of $\posC\op$ with $\poselc'' \posCleq \poselc$.
    The relation $\monrelM$ relates $\maketupconcat{\posela, \poselc}$ to $\maketupconcat{\poselb, \poselc'}$.
    The braiding brings $\poselc'$ and $\poselc''$ into the right order, and then the evaluation $\ev_{\posC}$ asks that $\poselc' \posCleq \poselc''$.
    An element $\poselc''$ with $\poselc' \posCleq \poselc'' \posCleq \poselc$ exists precisely when $\poselc' \posCleq \poselc$.
    Using that $\monrelM$ and the identity morphisms are closed downward in the source and upward in the target, one finds that $\Tr_{\posA,\posB}^{\posC}(\monrelM)$ relates $\posela$ to $\poselb$ precisely when there exist elements $\poselc$ and $\poselc'$ of $\posC$ such that
    \begin{equation}
        \tup{\maketupconcat{\posela, \poselc}, \maketupconcat{\poselb, \poselc'}} \setin \monrelM
        \qquad \text{and} \qquad
        \poselc' \posCleq \poselc.
    \end{equation}
    This is precisely the second description \cref{eq:mRelL-feedback-leq} of the trace for monotone relations of \cref{sec:trace-mRelL}.
    So that trace is the canonical trace of $\mRelL$.
    We also see where the inequality comes from which we met in codesign (\cref{sec:feedback-codesign}), where what is required has to be at most what is provided: it is built into the evaluation and the coevaluation.
\end{example}

\subsection{Canonical trace for linear relations}

We now show that the symmetric strict monoidal category of real linear relations is \SY{compact closed}, and therefore admits a canonical trace.

Recall this category from \cref{rem:LinRel-strict-monoidal} and \cref{sec:trace-LinRel}: its objects are the coordinate spaces $\reals^n$, and a morphism from $\reals^k$ to $\reals^m$ is a linear relation, that is, a linear subspace $\relA \subseteq \reals^{k+m}$, whose elements we write as pairs $\tup{\mat{u}, \mat{y}}$ with $\mat{u} \setin \reals^k$ and $\mat{y} \setin \reals^m$.
The monoidal product sums the dimensions, $\reals^k \mtimescat \reals^l = \reals^{k+l}$, and the monoidal unit is $\reals^0$, whose only element is the vector with no entries.
In this subsection we denote this category by $\stylecat{LinRel}_\reals$.

\subsubsection{Duality in~$\stylecat{LinRel}_\reals$}

We claim that every object in~$\stylecat{LinRel}_\reals$ is \SY{right dualizable}, and therefore the category is \SY{compact closed}.
In fact, the dual of any object is the object itself: for any $n \setin \natnumbers$, we set
\begin{equation}
    \dualof{(\reals^n)} = \reals^n.
\end{equation}
The evaluation and coevaluation are defined as follows.
The evaluation
\begin{equation}
    \ev_{\reals^n} \colon \reals^n \mtimescat \reals^n \mto \reals^0
\end{equation}
is the linear relation which relates a vector $\mat{v} \setin \reals^n$, stacked on top of a vector $\mat{w} \setin \reals^n$, to the only element of $\reals^0$ precisely when
\begin{equation}
    \mat{v} + \mat{w} = \mat{0}.
\end{equation}
In other words, as a subspace of $\reals^{2n+0} = \reals^{2n}$, it consists of the vectors with $\mat{w} = -\mat{v}$.

The coevaluation
\begin{equation}
    \coev_{\reals^n} \colon \reals^0 \mto \reals^n \mtimescat \reals^n
\end{equation}
is the linear relation which relates the only element of $\reals^0$ to a vector $\mat{v} \setin \reals^n$, stacked on top of a vector $\mat{w} \setin \reals^n$, precisely when $\mat{v} + \mat{w} = \mat{0}$.

\begin{remark}
    One can verify that~$\ev_{\reals^n}$ and~$\coev_{\reals^n}$ are indeed linear subspaces, and that the snake equations (\cref{eq:snake-eq-1,eq:snack-eq-2}) are satisfied.
    The key point is that the condition~$\mat{v} + \mat{w} = \mat{0}$ defines a linear subspace.
    For example, the left-hand side of \cref{eq:snake-eq-1} relates a vector $\mat{u}$ to a vector $\mat{y}$ precisely when there is a vector $\mat{w}$ with $\mat{y} + \mat{w} = \mat{0}$ and $\mat{w} + \mat{u} = \mat{0}$, that is, when $\mat{y} = \mat{u}$; this is the identity relation.
\end{remark}

\subsubsection{The canonical trace formula for linear relations}

Since~$\stylecat{LinRel}_\reals$ is \SY{compact closed}, the canonical trace (\cref{def:trace_gen_endo}) is defined.
Given a linear relation~$\relA \colon \reals^m \mtimescat \reals^k \mto \reals^n \mtimescat \reals^k$, that is, a linear subspace
\begin{equation}
    \relA \subseteq \reals^{(m+k)+(n+k)},
\end{equation}
the trace over~$\reals^k$ is the linear relation
\begin{equation}
    \Tr_{\reals^m, \reals^n}^{\reals^k}(\relA) \colon \reals^m \mto \reals^n
\end{equation}
given by
\begin{equation}
    \label{eq:canonical-trace-linrel}
    \Tr_{\reals^m, \reals^n}^{\reals^k}(\relA) = \makeset{ \tup{\mat{u}, \mat{y}} \mid \exists \ \mat{x} \setin \reals^k \colon
        \tup{\begin{bmatrix} \mat{u} \\ \mat{x} \end{bmatrix}, \begin{bmatrix} \mat{y} \\ \mat{x} \end{bmatrix}} \setin \relA },
\end{equation}
where $\mat{u} \setin \reals^m$ and $\mat{y} \setin \reals^n$.
This is precisely the trace for linear relations which we defined directly in \cref{eq:LinRel-trace}.

In other words, the trace identifies the ``input copy'' and the ``output copy'' of~$\reals^k$: one asks for the existence of a vector~$\mat{x} \setin \reals^k$ that serves simultaneously as the~$\reals^k$-component of both the input and the output.

\begin{remark}
    Comparing \cref{eq:canonical-trace-linrel} with the trace for~$\RelL$ (\cref{def:RelL-feedback,ex:RelDualsTrace}), we see the same pattern: the trace is obtained by ``identifying the feedback wire'', i.e., requiring that the component on the input side equals (in~$\RelL$) or sums to zero with (in the evaluation for~$\stylecat{LinRel}_\reals$) the component on the output side.
    In the present case, the condition~$\mat{v} + \mat{w} = \mat{0}$ in the (co)evaluation means that the feedback wire carries~$\mat{x}$ on one end and~$-\mat{x}$ on the other; after composing, this yields the condition that the same~$\mat{x}$ appears in both slots of~$\relA$.
\end{remark}
